\SourceLeaf{358}%
Earth as a whole. This allows transformations of the formulas that make them much more flexible and convenient. Indeed, although the transformed formulas are really only approximations, they can give as much or even more accuracy than the strictly exact formulas. This is something one. will not find paradoxical when one considers that the exact formulas must in any case be applied with the aid of trigonometrical tables, whose accuracy is finite and limited by the number of decimal digits. Among the various formulas given for this purpose in the first memoir, the combination at its close is especially noteworthy. It expresses the relationship between the six quantities in the most convenient form for calculation and gives more than enough accuracy even for the largest triangles that can actually be measured. This was bound to create a desire for analogous formulas that apply directly to the ellipsoid, and developing them is the main subject of the present \textit{second} memoir.

The formulas for the sphere mentioned above were found using entirely elementary propositions. Finding their counterparts on the ellipsoid, however, requires a series of fairly intricate operations. It is therefore undoubtedly welcome to be shown more than one way of reaching the same goal. The author carried out all these investigations for his own use more than thirty years ago and had so far found no particular occasion to publish them. In the present memoir he now gives \textit{two} entirely different derivations that lead in the end to exactly the same results. One of them is rooted in the theory of conformally mapping the ellipsoid onto the sphere. In this respect the second memoir also follows on from the first. Yet the two are to be regarded as entirely independent of each other insofar as one remains free to carry out the geodetic calculations solely by the method taught in either the first or the second memoir.

The problem of passing from the geographical position of a point on the spheroid to that of another point joined to it by a %
